\section{Ablauf des Live-Konzerts}
Das gesamte Konzert findet in einer Geschwindigkeit von 128 Schlägen pro Minute statt. Dies wird erreicht indem zu Beginn einmal der Befehl $setcps (128/60/4)$ ausgeführt wird. Das Konzert beginnt mit dem Klang von zwei Klangschalen des gleichen Tones. Anschließend wird der Klang einer der beiden Klangschalen mit dem $fshift$ Effekt von TidalCycles modifiziert. Durch die leichte Frequenzverschiebung einer der beiden Töne entsteht eine Schwebung (siehe Abschnitt~\ref{sec:schwebung}). Der Effekt ist dabei so eingestellt, das eine Frequenzverschiebung von $k \cdot \frac{128}{60}$ mit $k \in \mathbb{N}$ stattfindet. Dadurch wird eine Schwebungsfrequenz erzeugt, die ein ganzzahliges Vielfaches der Taktrate beträgt. Die Wellenbäuche bzw. -Knoten der Schwebung passen damit zum Takt der Musik. Die Umsetzung in TidalCycles ist in Listing~\ref{lst:tidal-schwebung} zu sehen. Die Funktion $gong$ erzeugt hierbei zwei Töne von welchen einer mit dem $fshift$ Effekt modifiziert wird. Der Parameter $shift$ entspricht dabei dem oben genannten $k$ und kann beliebig verändert werden, um die gewünschte Schwebungsfrequenz zu erzeugen.

\begin{lstlisting}[language=Haskell, caption={Erzeugung von $supergong$-Tönen mit Parametrisierung der Schwebungsfrequenz.}, label={lst:tidal-schwebung}]
let bpm = 128
setcps (bpm/60/4)

gong shift = stack [ slow 2 $ n (scale "hexSus" "f3") # s "supergong"
                   , slow 2 $ n (scale "hexSus" "f3") # s "supergong" # fshift (shift * bpm/60)
                   ]

d1 $ gong 2 # gain 0.6
\end{lstlisting}

Nun werden geschlossene Highhats hinzugefügt, welche mit einem Tiefpassfilter moduliert werden. Die Filterfrequenz wird hierbei mit einem Sinus der halben Schwebungsfrequenz geändert.

Nun wird ein eigens konstruierter Bass zugeschalten, dieser wurde mithilfe des Software-Synthesizers, der in Abschnitt~\ref{sec:softwaresynth} beschrieben wird, modular zusammengebaut und besteht aus zwei Sägezahnschwingungen und besitzt verschiedene Parameter welche über Drehknöpfe des MIDI-Keyboards, angepasst werden können (vgl. Abschnitt~\ref{sec:MIDI}).

%TODO: Erklärung aller Parameter ergänzen / genauer formulieren (vor allem Acidität)
Die vier einstellbaren Parameter sind eine Frequenz-Verschiebung einer der beiden Sägezahnschwingungen (ähnlich wie bei Klangschale um $k \cdot \frac{128}{60}$) um eine gezielte Schwebung zu erzeugen, ein Tiefpassfilter mithilfe welchem der Klang zwischen Sägezahn und harmonischer Schwingung verändert werden kann (vgl. Abschnitt \ref{sec:sawtooth}), ein Lautstärke-Regler und eine Einstellung mit welcher ein Filter der wiederum an den Attack-Wert der beiden Schwingungen gebunden ist eingestellt werden kann.

Der Bass wird zunächst mit einem Filterwert nahe der Grundfrequenz des Tons und einer Schwebung gleich der der Klangschalen verwendet. Nun wird zunächst ein Filter auf die Klangschalen angewendet bevor diese ganz entfernt werden.

Verschiedene Schlag-Elemente werden zugeschaltet, sowie der \textit{Korg Volca Keys}-Synthesizer, dieser spielt zunächst jedoch nur einzelne Töne zu Beginn jedes zweiten Taktes und keine Melodie. Durch Einstellen eines hohen Delay-Feedbacks und Veränderung der Zeit-Komponente wird eine Art Tremolo mit schwankender Tonhöhe erzeugt (siehe Abschnitt \ref{sec:delay}). \\

%TODO: warum hexSus und warum diese Töne, bringt der lpf wirklich was?? %
Ein Pad wird aus drei zyklisch wechselnden Akkorden innerhalb der hexSus-Skala erstellt (Vergleich Listing~\ref{lst:tidal-pad}). Das Pad erzeugt Hintergrundtextur und fügt gleichzeitig eine melodische Ebene hinzu.
Eine lange Attack-Zeit von drei Sekunden sowie eine Release-Zeit von fünf Sekunden sorgt dafür, dass die Akkorde ohne hörbaren Einsatzpunkt einschwingen und weich ineinander übergehen. In Kombination mit \texttt{legato 2} überlappen sich aufeinanderfolgende Akkorde zeitlich, sodass der perkussive Charakter eines klassischen Akkordanschlags vermieden wird. Ein stark ausgeprägter Reverb (\texttt{room 0.9}, \texttt{size 0.95}) verleiht dem Pad abschließend die räumliche Tiefe für den sphärischen Gesamtklang.

\begin{lstlisting}[language=Haskell, caption={Erzeugung Hintergrund Pad}, label={lst:tidal-pad}]
d3 $ slow 2 $ n ( scale "hexSus" "<[-2, 1, 5] [1, 3, 4, 7] c'min>") # s "superzow"
  # attack 3                            
  # release 5
  # sustain 1
  # legato 2
  # lpf (slow 20 $ range 400 "2000" $ sine)
  # room 0.9 # size 0.95
  # gain 0.6
\end{lstlisting}

\begin{enumerate}
    \item Einstellung Synthi "Acid" (rauer Ton durch hohe LFO Frequenz)
    \item Random Töne Synthi
    \item Abbau und Rückkehr zu Schwebung
\end{enumerate}